\documentclass[UTF8,fontset=windows,11pt]{ctexart}
\usepackage[a4paper,margin=1.85cm,headheight=24pt,footskip=24pt]{geometry}
\usepackage{xcolor}
\usepackage{graphicx}
\usepackage{enumitem}
\usepackage{tcolorbox}
\usepackage{fancyhdr}
\usepackage{lastpage}
\usepackage{titlesec}
\usepackage{hyperref}
\tcbuselibrary{skins,breakable}

\definecolor{examnavy}{HTML}{1E3A5F}
\definecolor{examblue}{HTML}{2563A7}
\definecolor{examred}{HTML}{C2410C}
\definecolor{examgray}{HTML}{4B5563}
\definecolor{examline}{HTML}{D7DEE8}
\definecolor{exambg}{HTML}{F8FAFC}
\definecolor{choicepink}{HTML}{DB2777}
\definecolor{answergreen}{HTML}{15803D}

\hypersetup{colorlinks=true,linkcolor=examnavy,urlcolor=examblue}
\setlength{\parindent}{0pt}
\setlength{\parskip}{0.28em}
\linespread{1.08}

\pagestyle{fancy}
\fancyhf{}
\lhead{\small\textcolor{examnavy}{施工章节测试}}
\rhead{\small\textcolor{examred}{复习指南对应题｜答案跟题}}
\cfoot{\small\textcolor{examgray}{第 \thepage\ 页 / 共 \pageref{LastPage} 页}}
\renewcommand{\headrulewidth}{0.4pt}
\renewcommand{\headrule}{\hbox to\headwidth{\color{examline}\leaders\hrule height \headrulewidth\hfill}}

\titleformat{\section}
  {\Large\bfseries\color{examnavy}}
  {\thesection}
  {0.8em}
  {}
\titlespacing*{\section}{0pt}{1.1em}{0.55em}

\newcommand{\chaptermeta}[1]{%
  \vspace{-0.25em}
  {\small\textcolor{examgray}{#1}}\par
  \vspace{0.35em}
  {\color{examline}\rule{\linewidth}{0.55pt}}\vspace{0.25em}
}

\newcommand{\choice}[1]{%
  \tcbox[
    on line,
    colback=white,
    colframe=choicepink,
    boxrule=0.55pt,
    arc=0.8mm,
    left=1.35mm,
    right=1.35mm,
    top=0.15mm,
    bottom=0.15mm
  ]{\textbf{\textcolor{choicepink}{#1}}}%
}

\newcommand{\questionstem}[1]{%
  \textbf{\textcolor{examnavy}{#1}}\par\vspace{0.35em}
}

\newcommand{\answerarea}[1]{%
  \vspace{0.2em}
  \begin{tcolorbox}[
    enhanced,
    colback=exambg,
    colframe=examline,
    boxrule=0.45pt,
    arc=0.8mm,
    height=#1,
    left=2mm,
    right=2mm,
    top=1mm,
    bottom=1mm
  ]
  \textcolor{examline}{作答区}
  \end{tcolorbox}
}

\newtcolorbox{questionbox}[2]{
  enhanced,
  breakable,
  colback=white,
  colframe=examline,
  boxrule=0.55pt,
  arc=1mm,
  left=3mm,
  right=3mm,
  top=2.4mm,
  bottom=2mm,
  before skip=0.65em,
  after skip=0.72em,
  borderline west={1.8pt}{0pt}{examblue},
  title={\textcolor{examnavy}{#1}\hfill\textcolor{examgray}{#2}},
  coltitle=examnavy,
  colbacktitle=exambg,
  fonttitle=\bfseries,
  boxed title style={
    colback=exambg,
    colframe=examline,
    boxrule=0.35pt,
    arc=0.8mm
  },
  attach boxed title to top left={xshift=2mm,yshift=-2mm},
  top=6mm
}


\newcommand{\guidebox}[1]{%
  \begin{tcolorbox}[
    enhanced,
    breakable,
    colback=examblue!5,
    colframe=examline,
    boxrule=0.4pt,
    arc=0.8mm,
    left=2mm,
    right=2mm,
    top=0.9mm,
    bottom=0.9mm
  ]
  {\small\textbf{\textcolor{examblue}{对应复习指南：}}#1}
  \end{tcolorbox}
}

\newcommand{\inlineanswer}[1]{%
  \vspace{0.25em}
  \begin{tcolorbox}[
    enhanced,
    breakable,
    colback=answergreen!6,
    colframe=answergreen,
    boxrule=0.55pt,
    arc=0.8mm,
    left=2.5mm,
    right=2.5mm,
    top=1.2mm,
    bottom=1.2mm,
    borderline west={1.6pt}{0pt}{answergreen}
  ]
  \textbf{\textcolor{answergreen}{答案：}}#1
  \end{tcolorbox}
}

\begin{document}



\begin{titlepage}
\vspace*{1.2cm}
{\Huge\bfseries\textcolor{examnavy}{施工章节测试}}\par
\vspace{0.35em}
{\LARGE\bfseries\textcolor{examred}{复习指南对应题 去重答案跟题版}}\par
\vspace{0.8em}
{\color{examline}\rule{0.82\linewidth}{0.8pt}}\par
\vspace{1.1em}
{\large 仅保留参考 PDF 中与复习指南对应的题目；同一知识点或题干高度相似的题已去重。}\par
\vspace{2.2em}
\begin{tcolorbox}[
  enhanced,
  width=0.78\linewidth,
  colback=exambg,
  colframe=examline,
  boxrule=0.6pt,
  arc=1mm,
  left=4mm,
  right=4mm,
  top=3mm,
  bottom=3mm
]
\Large
\textcolor{examnavy}{参考题数}\quad\textbf{101}\qquad
\textcolor{examnavy}{去重后}\quad\textbf{21}\qquad
\textcolor{examnavy}{删除相似}\quad\textbf{80}\par
\vspace{0.45em}
\normalsize\textcolor{examgray}{顺序版；每题后直接附答案，不生成乱序版。}
\end{tcolorbox}
\vfill
{\small\textcolor{examgray}{依据《施工章节测试\_复习指南对应题\_刷题版4.pdf》重新整理}}
\end{titlepage}
\tableofcontents
\newpage


\section{第一章土方工程章节测试}

\chaptermeta{本章保留 5 题；题号括号中标注原全卷题号。}

\begin{questionbox}{第 1 题（原第 1 题）}{单选题}
\guidebox{土方工程：填土压实质量控制指标——土的干密度}
\questionstem{作为检验填土压实质量控制指标的是（）。}
\begin{enumerate}[label=\protect\choice{\Alph*}, leftmargin=*, itemsep=0.32em, topsep=0.2em]
\item 土的干密度
\item 土的密实度
\item 土的压缩比
\item 土的可松性
\end{enumerate}
\inlineanswer{A}
\end{questionbox}

\begin{questionbox}{第 2 题（原第 7 题）}{多选题}
\guidebox{土方工程：土方施工机械的选择因素}
\questionstem{土方机械的选择，一般应考虑的因素有（）。}
\begin{enumerate}[label=\protect\choice{\Alph*}, leftmargin=*, itemsep=0.32em, topsep=0.2em]
\item 土方工程的范围大小
\item 施工合同的要求
\item 不同施工机械的组合
\item 土的含水量
\item 施工合同工期要求
\end{enumerate}
\inlineanswer{ACDE}
\end{questionbox}

\begin{questionbox}{第 3 题（原第 10 题）}{单选题}
\guidebox{土方工程：基坑挖土方案——放坡、中心岛、盆式与逆作法}
\questionstem{当地质条件和场地条件许可时，开挖深度不大的基坑可采用的开挖方案是（）。}
\begin{enumerate}[label=\protect\choice{\Alph*}, leftmargin=*, itemsep=0.32em, topsep=0.2em]
\item 放坡挖土
\item 中心岛挖土
\item 盆式挖土
\item 逆作法挖土
\end{enumerate}
\inlineanswer{A}
\end{questionbox}

\begin{questionbox}{第 4 题（原第 17 题）}{多选题}
\guidebox{土方工程：回填土料、分层填筑、分层压实与分层检测}
\questionstem{填方土料应符合设计要求，一般不能作为回填土的有（）。}
\begin{enumerate}[label=\protect\choice{\Alph*}, leftmargin=*, itemsep=0.32em, topsep=0.2em]
\item 淤泥
\item 淤泥质土
\item 有机质为8\%的土
\item 黏性土
\item 砂砾坚土
\end{enumerate}
\inlineanswer{ABC}
\end{questionbox}

\begin{questionbox}{第 5 题（原第 24 题）}{多选题}
\guidebox{土方工程：基坑验槽的组织、参加单位、方法与检查内容}
\questionstem{可以组织基坑验槽的人员有（）。}
\begin{enumerate}[label=\protect\choice{\Alph*}, leftmargin=*, itemsep=0.32em, topsep=0.2em]
\item 设计单位负责人
\item 总监理工程师
\item 施工单位负责人
\item 建设单位项目负责人
\item 勘察单位负责人
\end{enumerate}
\inlineanswer{BD}
\end{questionbox}

\section{第二章桩基础工程章节测试}

\chaptermeta{本章保留 1 题；题号括号中标注原全卷题号。}

\begin{questionbox}{第 6 题（原第 37 题）}{单选题}
\guidebox{桩基础工程：静压桩终止沉桩控制原则}
\questionstem{下列关于静压桩终止沉桩标准的说法，错误的是（）。}
\begin{enumerate}[label=\protect\choice{\Alph*}, leftmargin=*, itemsep=0.32em, topsep=0.2em]
\item 静压桩应以标高为主，压力为辅
\item 摩擦桩应按桩顶标高控制
\item 端承桩应以终压力控制为主，桩顶标高控制为辅
\item 端承摩擦桩，应以终压力控制为主，桩顶标高控制为辅
\end{enumerate}
\inlineanswer{D}
\end{questionbox}

\section{第三章基坑工程章节测试}

\chaptermeta{本章保留 2 题；题号括号中标注原全卷题号。}

\begin{questionbox}{第 7 题（原第 63 题）}{多选题}
\guidebox{基坑工程：集水井、降水井、减压井及承压水突涌控制}
\questionstem{可以起到防止深基坑坑底突涌的措施有（）。}
\begin{enumerate}[label=\protect\choice{\Alph*}, leftmargin=*, itemsep=0.32em, topsep=0.2em]
\item 集水明排
\item 钻孔减压
\item 井点降水
\item 井点回灌
\item 水平封底隔渗
\end{enumerate}
\inlineanswer{BE}
\end{questionbox}

\begin{questionbox}{第 8 题（原第 67 题）}{单选题}
\guidebox{基坑工程：井点回灌技术的作用与布置}
\questionstem{工程基坑开挖常用井点回灌技术的主要目的是（）。}
\begin{enumerate}[label=\protect\choice{\Alph*}, leftmargin=*, itemsep=0.32em, topsep=0.2em]
\item 避免坑底土体回弹
\item 避免坑底出现管涌
\item 减少排水设施，降低施工成本
\item 防止降水井点对井点周围建（构）筑物、地下管线的影响
\end{enumerate}
\inlineanswer{D}
\end{questionbox}

\section{第四章混凝土结构工程章节测试}

\chaptermeta{本章保留 6 题；题号括号中标注原全卷题号。}

\begin{questionbox}{第 9 题（原第 77 题）}{单选题}
\guidebox{混凝土结构工程：成型钢筋进场复验项目}
\questionstem{成型钢筋在进场时无需复验的项目是（）。}
\begin{enumerate}[label=\protect\choice{\Alph*}, leftmargin=*, itemsep=0.32em, topsep=0.2em]
\item 抗拉强度
\item 弯曲性能
\item 伸长率
\item 重量偏差
\end{enumerate}
\inlineanswer{B}
\end{questionbox}

\begin{questionbox}{第 10 题（原第 109 题）}{单选题}
\guidebox{混凝土结构工程：底模及支架拆除的同条件养护强度要求}
\questionstem{某跨度8m的混凝土楼板，设计强度等级为C30。模板采用快拆支架体系，支架立杆间距为2m，拆模时混凝土的最低强度是（）MPa。}
\begin{enumerate}[label=\protect\choice{\Alph*}, leftmargin=*, itemsep=0.32em, topsep=0.2em]
\item 15
\item 22.5
\item 25.5
\item 30
\end{enumerate}
\inlineanswer{A}
\end{questionbox}

\begin{questionbox}{第 11 题（原第 121 题）}{单选题}
\guidebox{混凝土结构工程：大体积混凝土防裂、材料选择与温度控制}
\questionstem{配制厚大体积的普通混凝土不宜选用（）水泥。}
\begin{enumerate}[label=\protect\choice{\Alph*}, leftmargin=*, itemsep=0.32em, topsep=0.2em]
\item 矿渣
\item 粉煤
\item 复合
\item 硅酸盐
\end{enumerate}
\inlineanswer{D}
\end{questionbox}

\begin{questionbox}{第 12 题（原第 127 题）}{多选题}
\guidebox{混凝土与砌筑工程：坍落度、维勃稠度、砂浆稠度及流动性}
\questionstem{下列关于砂浆、混凝土流动性说法中正确的有（）。}
\begin{enumerate}[label=\protect\choice{\Alph*}, leftmargin=*, itemsep=0.32em, topsep=0.2em]
\item 稠度测定仪测得的砂浆稠度值越小，流动性越大
\item 稠度测定仪测得的砂浆稠度值越大，流动性越大
\item 混凝土坍落度与坍落扩展度越小，流动性越大
\item 混凝土坍落度与坍落扩展度越大，流动性越大
\item 维勃稠度仪测得的混凝土稠度值越小，流动性越大
\end{enumerate}
\inlineanswer{BDE}
\end{questionbox}

\begin{questionbox}{第 13 题（原第 138 题）}{单选题}
\guidebox{混凝土结构工程：减水剂、引气剂、早强剂与速凝剂}
\questionstem{下列混凝土外加剂中，不能显著改善混凝土拌合物流变性能的是（）。}
\begin{enumerate}[label=\protect\choice{\Alph*}, leftmargin=*, itemsep=0.32em, topsep=0.2em]
\item 减水剂
\item 引气剂
\item 膨胀剂
\item 泵送剂
\end{enumerate}
\inlineanswer{C}
\end{questionbox}

\begin{questionbox}{第 14 题（原第 146 题）}{单选题}
\guidebox{混凝土结构工程：混凝土分层浇筑与振捣要求}
\questionstem{下列关于混凝土浇筑振捣的说法，正确的是（）。}
\begin{enumerate}[label=\protect\choice{\Alph*}, leftmargin=*, itemsep=0.32em, topsep=0.2em]
\item 混凝土不宜分层浇筑，分层振动
\item 插入式振捣棒振捣混凝土时，应慢插快拔
\item 插入式振捣棒移动间距不宜大于振捣棒作用半径的1.4倍
\item 振捣棒插入下层混凝土内的深度应不小于30mm
\end{enumerate}
\inlineanswer{C}
\end{questionbox}

\section{第五章预应力混凝土工程章节测试}

\chaptermeta{本章保留 1 题；题号括号中标注原全卷题号。}

\begin{questionbox}{第 15 题（原第 162 题）}{单选题}
\guidebox{预应力混凝土工程：无粘结预应力施工特点与工序}
\questionstem{无粘结后张法预应力混凝土中的预应力，是靠（）传递给混凝土。}
\begin{enumerate}[label=\protect\choice{\Alph*}, leftmargin=*, itemsep=0.32em, topsep=0.2em]
\item 锚具
\item 台座
\item 张拉千斤顶
\item 混凝土与预应力筋之间的粘结力
\end{enumerate}
\inlineanswer{A}
\end{questionbox}

\section{第六章砌筑工程章节测试}

\chaptermeta{本章保留 2 题；题号括号中标注原全卷题号。}

\begin{questionbox}{第 16 题（原第 180 题）}{多选题}
\guidebox{砌筑工程：“三一”砌筑法}
\questionstem{砖砌体“三一”砌筑法的具体含义是指（）。}
\begin{enumerate}[label=\protect\choice{\Alph*}, leftmargin=*, itemsep=0.32em, topsep=0.2em]
\item 一个人
\item 一铲灰
\item 一块砖
\item 一挤揉
\item 一勾缝
\end{enumerate}
\inlineanswer{BCD}
\end{questionbox}

\begin{questionbox}{第 17 题（原第 195 题）}{多选题}
\guidebox{砌筑工程：普通混凝土小型空心砌块施工要求}
\questionstem{关于普通混凝土小砌块施工说法，正确的是（）。}
\begin{enumerate}[label=\protect\choice{\Alph*}, leftmargin=*, itemsep=0.32em, topsep=0.2em]
\item 在施工前先浇水湿透
\item 清除表面污物
\item 底面朝下正砌于墙上
\item 底面朝上反砌于墙上
\item 小砌块在使用时的龄期已到
\end{enumerate}
\inlineanswer{BD}
\end{questionbox}

\section{第七章脚手架工程章节测试}

\chaptermeta{本章保留 1 题；题号括号中标注原全卷题号。}

\begin{questionbox}{第 18 题（原第 202 题）}{单选题}
\guidebox{危大工程管理：专项施工方案专家论证范围}
\questionstem{下列分部分项工程中，其专项施工方案不需要进行专家论证的有（）。}
\begin{enumerate}[label=\protect\choice{\Alph*}, leftmargin=*, itemsep=0.32em, topsep=0.2em]
\item 架体高度23m的悬挑脚手架
\item 开挖深度10m的人工挖孔桩
\item 爬升高度80m的爬模
\item 埋深10m的地下暗挖
\end{enumerate}
\inlineanswer{B}
\end{questionbox}

\section{第九章防水工程章节测试}

\chaptermeta{本章保留 1 题；题号括号中标注原全卷题号。}

\begin{questionbox}{第 19 题（原第 213 题）}{多选题}
\guidebox{防水工程：屋面卷材防水施工顺序、铺贴与搭接缝}
\questionstem{屋面防水施工基本要求正确的是（）。}
\begin{enumerate}[label=\protect\choice{\Alph*}, leftmargin=*, itemsep=0.32em, topsep=0.2em]
\item 以排为主，以防为辅
\item 上下层卷材不得相互垂直铺贴
\item 屋面卷材防水施工时，由高向低铺贴
\item 天沟卷材施工时，宜顺天沟方向铺贴
\item 立面或大坡面贴卷材应采用满粘法
\end{enumerate}
\inlineanswer{BDE}
\end{questionbox}

\section{第十章流水施工原理章节测试}

\chaptermeta{本章保留 1 题；题号括号中标注原全卷题号。}

\begin{questionbox}{第 20 题（原第 219 题）}{单选题}
\guidebox{逻辑思维内容指南1：流水施工原理}
\questionstem{关于横道图进度计划特点的说法，正确的是（）。}
\begin{enumerate}[label=\protect\choice{\Alph*}, leftmargin=*, itemsep=0.32em, topsep=0.2em]
\item 可以识别计划的关键工作
\item 不能表达工作逻辑关系
\item 调整计划的工作量较大
\item 可以计算工作时差
\end{enumerate}
\inlineanswer{C}
\end{questionbox}

\section{第十一章网络计划技术章节测试}

\chaptermeta{本章保留 1 题；题号括号中标注原全卷题号。}

\begin{questionbox}{第 21 题（原第 243 题）}{单选题}
\guidebox{逻辑思维内容指南2：单目标双代号网络计划六大时间参数}
\questionstem{工程网络计划中，工作的最迟开始时间是指在不影响（）的前提下，必须开始的最迟时刻。}
\begin{enumerate}[label=\protect\choice{\Alph*}, leftmargin=*, itemsep=0.32em, topsep=0.2em]
\item 紧后工作最早开始
\item 紧前工作最迟开始
\item 整个任务按期完成
\item 所有后续工作机动时间
\end{enumerate}
\inlineanswer{C}
\end{questionbox}

\end{document}
